\documentclass{ximera}


% MATH -----------------------------------------------------------
\newcommand{\nats}{\mathbb N}
\newcommand{\ints}{\mathbb Z}
\newcommand{\rats}{\mathbb Q}
\newcommand{\reals}{\mathbb R}
\newcommand{\complex}{\mathbb C}
\newcommand{\powerset}{\mathscr P}

%-------------------------------------------------------------

\newcounter{NumberInTable}
\newcommand{\LTNUM}{\stepcounter{NumberInTable}{(\theNumberInTable)}}

\newcommand{\Laplace}[1]{\ensuremath{\mathcal{L}{\left[#1\right]}}}
\newcommand{\InvLap}[1]{\ensuremath{\mathcal{L}^{-1}{\left[#1\right]}}}


\begin{document}
\begin{abstract}
 \end{abstract}

    \title{Inverse Laplace Transforms}
       \maketitle

 For the types of functions we encounter in this course, the Laplace transform is $1-1$.\\

%%%% SO IT MAKES SENSE TO CONSIDER THE INVERSE TRANSFORM

 Suppose $F(s)=\mathcal{L}(f(t))$.  For example, $\frac{6}{s^4}=\mathcal{L}(t^3)$.  Then we call $f(t)$ the \underline{inverse Laplace transform} of $F(s)$.   In our example, $t^3$ is the inverse Laplace transform of $\frac{6}{s^4}$.   Notation:  $f=\mathcal{L}^{-1}(F)$.\\


 While there is a formula for computing $\mathcal{L}^{-1}(F)$ we will not use it as it requires the theory of functions of a complex variable (complex analysis).   We can however, make use the table (see appendix) to find inverse Laplace transforms.\\


 Suppose we need $\mathcal{L}^{-1}\left(\dfrac{s}{s^{2}+4}\right)$.   By looking through the table of Laplace transforms we find $\mathcal{L}(\cos{(kt)})=\dfrac{s}{s^2+k^2}$.  So $\mathcal{L}^{-1}\left(\dfrac{s}{s^{2}+4}\right)=\cos{(2t)}$.\\


 Note:  Linearity also holds for the inverse Laplace transform! We leave the justification as exercise left for the reader.\\


 Let's find $\mathcal{L}^{-1}\left(\dfrac{2s+3}{s^{2}+4s+13}\right)$.\\


 $$\mathcal{L}^{-1}\left(\dfrac{2s+3}{s^{2}+4s+13}\right)=\mathcal{L}^{-1}\left(\dfrac{2(s+2)-1}{(s+2)^{2}+9}\right)=2\mathcal{L}^{-1}\left(\dfrac{s+2}{(s+2)^{2}+9}\right)-\dfrac{1}{3}\mathcal{L}^{-1}\left(\dfrac{3}{(s+2)^{2}+9}\right)$$

\medskip

 \hspace{0.5cm} $\displaystyle =2e^{-2t}\cos{(3t)}-\dfrac{1}{3}e^{-2t}\sin{(3t)}=e^{-2t}\left(2\cos{(3t)}-\dfrac{1}{3}\sin{(3t)}\right)$\\ \\


\newpage

 In the last example, the denominator is irreducible over the reals.  Our method changes when the denominator factors nicely:\\


 Let's find $\mathcal{L}^{-1}\left(\dfrac{2s-1}{s^{2}-5s+6}\right)$.\\

 Using a PFD (partial fraction decomposition)  we get \\$\displaystyle \dfrac{2s-1}{s^{2}-5s+6}= \dfrac{2s-1}{(s-2)(s-3)}=\dfrac{5}{s-3}-\dfrac{3}{s-2}$.\\


 Then $\mathcal{L}^{-1}\left(\dfrac{2s-1}{s^{2}-5s+6}\right)=5\mathcal{L}^{-1}\left(\dfrac{1}{s-3}\right)-3\mathcal{L}^{-1}\left(\dfrac{1}{s-2}\right)=5e^{3t}-3e^{2t}$.\\


 Let's find $\mathcal{L}^{-1}\left(\dfrac{s}{s^{2}-8s+16}\right)$.   We can setup a PFD as follows:  $$\dfrac{s}{s^{2}-8s+16}=\dfrac{s}{(s-4)^{2}}=\dfrac{A}{s-4}+\dfrac{4}{(s-4)^{2}}.$$

\bigskip

 Clearing fractions and equating common-degree coefficients, we get $A=1$.

\bigskip

 So $\displaystyle \mathcal{L}^{-1}\left(\dfrac{s}{s^{2}-8s+16}\right)=\mathcal{L}^{-1}\left(\dfrac{1}{s-4}\right)+4\mathcal{L}^{-1}\left(\dfrac{1}{(s-4)^{2}}\right)=e^{4t}+4te^{4t}=e^{4t}(1+4t)$.

\newpage


\begin{center}
\begin{LARGE}
APPENDIX: Table of Laplace Transforms
\end{LARGE}
\end{center}

\vspace{2ex}

\renewcommand{\arraystretch}{1.5}
\begin{center}
\begin{tabbing}
\hspace*{1.5 in}\=\hspace{1.5in}\= \kill
\hspace{0.5cm} $f(t)$ \> $\Laplace{f(t)}=F(s)$ \> \\ \\ \\
\LTNUM $\mbox{ }1$       			 \> $\dfrac{1}{s}$           \> \\ \\
\LTNUM $\mbox{ }e^{at}f(t)$	\> $F(s-a)$	\> \\  \\
\LTNUM $\mbox{ }\mathcal{U}(t-a)$ \> $\dfrac{e^{-as}}{s}$ \> \\ \\
\LTNUM $\mbox{ }f(t-a)\mathcal{U}(t-a)$ \> $e^{-as}F(s)$ \> \\ \\
\LTNUM $\mbox{ }\delta(t)$	\> 1 \>  \\ \\
\LTNUM $\mbox{ }\delta(t-t_0)$ \> $e^{-st_0}$ \> \\ \\
\LTNUM $\mbox{ }t^nf(t)$ 	\> $(-1)^n\dfrac{d^nF(s)}{ds^n}$  \> \\ \\
\LTNUM $\mbox{ }f'(t)$ 	\> $sF(s) - f(0)$ \> \\ \\
\LTNUM $\mbox{ }f^{n}(t)$ 	\> $s^nF(s) - \displaystyle \sum_{k=0}^{n-1} s^{(n-1-k)} f^{(k)}(0)$ \> \\  \\
\LTNUM $\displaystyle{\int_0^t f(x)g(t-x)dx}$ \> $F(s)G(s)$ \> \\ \\
\LTNUM $\mbox{ }t^n$ ($n=0,1,\dots$)     \> $\dfrac{n!}{s^{n+1}}$    \> \\ \\
\LTNUM $\mbox{ }t^x$ ($x\geq-1\in\mathbb{R}$)     \> $\dfrac{\Gamma(x+1)}{s^{x+1}}$ \> \\ \\
\LTNUM $\mbox{ }\sin kt$ 	\> $\dfrac{k}{s^2+k^2}$ \> \\ \\
\LTNUM $\mbox{ }\cos kt$ 	\> $\dfrac{s}{s^2+k^2}$ \> \\ \\
\LTNUM $\mbox{ }e^{at}$ 	\> $\dfrac{1}{s-a}$ 	\> 
 \end{tabbing}
 \end{center}

\newpage

\begin{center}
\begin{tabbing}
\hspace*{1.5in}\=\hspace{1.5in}\= \kill
\hspace{1cm} $f(t)$ \> $\Laplace{f(t)}=F(s)$ \> \\ \\ \\
\LTNUM $\mbox{ }\sinh kt$	\> $\dfrac{k}{s^2-k^2}$ \> \\ \\
\LTNUM $\mbox{ }\cosh kt$	\> $\dfrac{s}{s^2-k^2}$ \> \\ \\
\LTNUM $\mbox{ }\dfrac{e^{at}-e^{bt}}{a-b}$	\> $\dfrac{1}{(s-a)(s-b)}$ \> \\ 

\LTNUM $\mbox{ }\dfrac{ae^{at}-be^{bt}}{a-b}$	\> $\dfrac{s}{(s-a)(s-b)}$ \> \\ \\
\LTNUM $\mbox{ }te^{at}$	\> $\dfrac{1}{(s-a)^2}$	\> \\ \\
\LTNUM $\mbox{ }t^ne^{at}$	\> $\dfrac{n!}{(s-a)^{n+1}}$	\> \\ \\ 
\LTNUM $\mbox{ }e^{at}\sin kt$	\> $\dfrac{k}{(s-a)^2+k^2}$  \> \\ \\
\LTNUM $\mbox{ }e^{at}\cos kt$	\> $\dfrac{s-a}{(s-a)^2+k^2}$  \> \\ \\ 
\LTNUM $\mbox{ }e^{at}\sinh kt$	\> $\dfrac{k}{(s-a)^2-k^2}$  \> \\ \\
\LTNUM $\mbox{ }e^{at}\cosh kt$	\> $\dfrac{s-a}{(s-a)^2-k^2}$  \> \\ \\
\LTNUM $\mbox{ }t\sin kt$  	\> $\dfrac{2ks}{(s^2+k^2)^2}$ \> \\ \\
\LTNUM $\mbox{ }t\cos kt$  	\> $\dfrac{s^2-k^2}{(s^2+k^2)^2}$ \> \\  \\ 
\LTNUM $\mbox{ }t\sinh kt$  	\> $\dfrac{2ks}{(s^2-k^2)^2}$ \> \\ \\
\LTNUM $\mbox{ }t\cosh kt$  	\> $\dfrac{s^2-k^2}{(s^2-k^2)^2}$ \> \\ \\ 
\LTNUM $\mbox{ }\dfrac{\sin at}{t}$	\> $\arctan \dfrac{a}{s}$ \> 
 \end{tabbing}
 \end{center}

\renewcommand{\arraystretch}{1}




\end{document}